% !TEX root = task-6-8-1.tex
% Editable task bank. Append new BankTask blocks at the end of the document.
\documentclass[12pt,a4paper]{article}
\usepackage{physics-textbook}
\usepackage{graphicx,needspace}
\usetikzlibrary{decorations.pathmorphing}
\newif\ifShowSolutions
\ifdefined\TutorEdition
  \ShowSolutionstrue
\else
  \ShowSolutionsfalse % Change to \ShowSolutionstrue to display solutions.
\fi
\newcounter{banktask}
\newsavebox{\BankTaskBox}
\newcommand{\BankTaskContent}[3]{%
  \noindent\textbf{\thebanktask.}\quad #2\par
  \ifShowSolutions
    \Needspace{50mm}% Keep the solution heading with its opening text or diagram.
    #3\par
    \ifdefined\TutorEdition
      \ifcsname TutorSource#1\endcsname
        \nopagebreak\smallskip{\small\csname TutorSource#1\endcsname\par}%
      \fi
    \fi
  \fi
}
\newcommand{\BankTask}[3]{%
  \refstepcounter{banktask}\par\medskip
  \sbox{\BankTaskBox}{\begin{minipage}{\linewidth}%
    \BankTaskContent{#1}{#2}{#3}\end{minipage}}%
  \ifdim\dimexpr\ht\BankTaskBox+\dp\BankTaskBox\relax>\dimexpr\textheight-20mm\relax
    \Needspace{35mm}\BankTaskContent{#1}{#2}{#3}%
  \else
    \Needspace{\dimexpr\ht\BankTaskBox+\dp\BankTaskBox\relax}%
    \noindent\usebox{\BankTaskBox}\par
  \fi
  \medskip
}
\begin{document}
\BookSubsection{6.8.1}{Нуклонная модель ядра}


% Source sdamgia-3342: https://phys-ege.sdamgia.ru/problem?id=3342
\BankTask{sdamgia-3342}{
Каков заряд ядра ${}^{11}_{5}\mathrm{B}$ (в единицах элементарного заряда)?
}{
\smallskip\noindent\textbf{Решение}\par
Число протонов (заряд ядра в элементарных единицах заряда) записывается внизу перед наименованием элемента. Таким образом, заряд ядра ${}^{11}_{5}\mathrm{B}$ равен $5$.

Ответ: $5$.
}

% Source sdamgia-8013: https://phys-ege.sdamgia.ru/problem?id=8013
\BankTask{sdamgia-8013}{
Сколько протонов и сколько нейтронов содержится в ядре ${}^{60}_{27}\mathrm{Co}$?

\begin{center}\small\renewcommand{\arraystretch}{1.25}\begin{tabular}{|p{0.450\linewidth}|p{0.450\linewidth}|}\hline
Число протонов & Число нейтронов\\\hline
 & \\\hline
\end{tabular}\end{center}
}{
\smallskip\noindent\textbf{Решение}\par
В соответствии с обозначением изотопов в ядре ${}^{60}_{27}\mathrm{Co}$ $27$ протонов и $60 - 27 = 33$ нейтрона.

Ответ: $2733$.
}

% Source sdamgia-8869: https://phys-ege.sdamgia.ru/problem?id=8869
\BankTask{sdamgia-8869}{
Электронная оболочка электрически нейтрального атома криптона содержит $36$ электронов. Сколько нейтронов содержится в ядрах изотопов криптона-$78$ и криптона-$86$? 

В ответе запишите только числа, не разделяя их пробелом или другим знаком.

\begin{center}\small\renewcommand{\arraystretch}{1.25}\begin{tabular}{|p{0.450\linewidth}|p{0.450\linewidth}|}\hline
Число нейтронов в ядре

криптона-$78$ & Число нейтронов в ядре

криптона-$86$\\\hline
 & \\\hline
\end{tabular}\end{center}
}{
\smallskip\noindent\textbf{Решение}\par
В электрически нейтральном атоме количество электронов равно количеству протонов, значит, в атоме криптона содержится $36$ протонов. В ядре изотопа криптона-$78$ содержится $78 - 36 = 42$ нейтрона, В ядре изотопа криптона-$86$ — $86 - 36 = 50$ нейтронов.

Ответ: $4250$.
}

% Source sdamgia-10191: https://phys-ege.sdamgia.ru/problem?id=10191
\BankTask{sdamgia-10191}{
В ноябре $2016\,\text{г}$. в периодическую таблицу Менделеева был официально внесен новый химический элемент с порядковым номером $115$ — он получил название «московий» (обозначается $Mc$). Атомная масса наиболее стабильного изотопа этого элемента (из ныне известных) равна $289$. Сколько протонов и сколько нейтронов содержится в атомном ядре этого изотопа московия? В ответ запишите два числа без пробелов и запятых.

\begin{center}\small\renewcommand{\arraystretch}{1.25}\begin{tabular}{|p{0.450\linewidth}|p{0.450\linewidth}|}\hline
Число протонов & Число нейтронов\\\hline
 & \\\hline
\end{tabular}\end{center}
}{
\smallskip\noindent\textbf{Решение}\par
Номер элемента в таблице Менделеева показывает количество протонов, значит, у московия имеется $115$ протонов и $289 - 115 = 174$ нейтрона.

Ответ: $115174$.
}

% Source sdamgia-19672: https://phys-ege.sdamgia.ru/problem?id=19672
\BankTask{sdamgia-19672}{
В составе изотопа некоторого элемента число протонов в $1{,}5$ раза меньше числа нуклонов. Определите число протонов и число нуклонов в этом ядре, если известно, что его массовое число равно $6$.

\begin{center}\small\renewcommand{\arraystretch}{1.25}\begin{tabular}{|p{0.450\linewidth}|p{0.450\linewidth}|}\hline
Число протонов & Число нуклонов\\\hline
 & \\\hline
\end{tabular}\end{center}
}{
\smallskip\noindent\textbf{Решение}\par
$1$. Число протонов: массовое число / $1{,}5 = 4$ протона.

$2$. Число нуклонов = сумма протонов и нейтронов = массовое число = $6$.

Ответ: $46$.
}

% Source sdamgia-24108: https://phys-ege.sdamgia.ru/problem?id=24108
\BankTask{sdamgia-24108}{
В некотором атомном ядре число нейтронов в $2$ раза больше числа протонов. Массовое число этого ядра равно $36$. Определите для этого ядра число нейтронов и разность между числом нейтронов и зарядовым числом.

\begin{center}\small\renewcommand{\arraystretch}{1.25}\begin{tabular}{|p{0.450\linewidth}|p{0.450\linewidth}|}\hline
Число нейтронов & Разность между числом нейтронов

и зарядовым числом\\\hline
 & \\\hline
\end{tabular}\end{center}
}{
\smallskip\noindent\textbf{Решение}\par
Поскольку число нейтронов в $2$ раза больше, чем протонов, то число нуклонов в $3$ раза больше, чем протонов. Учитывая, что нуклонов в ядре $36$, находим число протонов — $12$, нейтронов — $24$. Тогда разница между этими числами равна $12$.

Ответ: $2412$.
}

% Source sdamgia-23242: https://phys-ege.sdamgia.ru/problem?id=23242
\BankTask{sdamgia-23242}{
Ядро атома золота имеет электрический заряд, который в $79$ раз превышает модуль заряда электрона. В настоящее время известны изотопы золота с числом нейтронов от $90$ до $126$. Определите для известных изотопов золота массовое число самого легкого изотопа, а также разность атомных масс самого тяжелого и самого легкого изотопов.

\begin{center}\small\renewcommand{\arraystretch}{1.25}\begin{tabular}{|p{0.450\linewidth}|p{0.450\linewidth}|}\hline
Массовое число самого легкого изотопа & Разность атомных масс самого тяжелого и самого легкого изотопов\\\hline
 & \\\hline
\end{tabular}\end{center}
}{
\smallskip\noindent\textbf{Решение}\par
Число протонов равно числу электронов. Из условия следует, что в ядре атома золота $79$ протонов. Массовое число равно сумме числа протонов и нейтронов. Значит, у легкого ядра массовое число равно $79 + 90 = 169$, у тяжелого $79 + 126 = 205$. Разность массовых чисел тяжелого и легкого ядер равна $205 - 169 = 36$.

Ответ: $16936$.
}

% Source sdamgia-24960: https://phys-ege.sdamgia.ru/problem?id=24960
\BankTask{sdamgia-24960}{
Атом изотопа бериллия ${}^{7}_{4}\mathrm{Be}$ подвергся двойной ионизации, в результате чего образовался положительный ион B$e^{2+}$. Сколько нейтронов содержит ядро этого атома, и сколько электронов осталось у иона?

\begin{center}\small\renewcommand{\arraystretch}{1.25}\begin{tabular}{|p{0.450\linewidth}|p{0.450\linewidth}|}\hline
Число нейтронов в ядре атома & Число электронов у иона\\\hline
 & \\\hline
\end{tabular}\end{center}
}{
\smallskip\noindent\textbf{Решение}\par
В атоме бериллия содержалось $4$ протона, $4$ электрона и $3$ нейтрона. При двойной ионизации число протонов и нейтронов не изменилось, электронов стало $2$. 

Ответ: $32$.
}

% Source sdamgia-9032: https://phys-ege.sdamgia.ru/problem?id=9032
\BankTask{sdamgia-9032}{
Сколько протонов и сколько нейтронов содержится в ядре ${}^{55}_{26}$Fe?

\begin{center}\small\renewcommand{\arraystretch}{1.25}\begin{tabular}{|p{0.450\linewidth}|p{0.450\linewidth}|}\hline
Число протонов & Число нейтронов\\\hline
 & \\\hline
\end{tabular}\end{center}
}{
\smallskip\noindent\textbf{Решение}\par
В ядре ${}^{55}_{26}$Fe содержится $26$ протонов и $55 - 26 = 29$ нейтронов.

Ответ: $2629$.
}

% Source sdamgia-6932: https://phys-ege.sdamgia.ru/problem?id=6932
\BankTask{sdamgia-6932}{
Сколько электронов вращается вокруг ядра атома ${}^{89}_{39}\mathrm{Y}?$
}{
\smallskip\noindent\textbf{Решение}\par
Заряд ядра атома иттрия равен $39$, а значит, в ядре $39$ протонов. В нейтральном атоме число протонов равно числу электронов. Следовательно, вокруг ядра вращается $39$ электронов.

Ответ: $39$.
}


% Source B00348: https://ege.fipi.ru/bank/questions.php?search=1&proj=BA1F39653304A5B041B656915DC36B38&qid=B00348
\BankTask{B00348}{
На рисунке представлен фрагмент Периодической системы элементов Д.И. Менделеева. Под названием каждого элемента приведены массовые числа его основных стабильных изотопов. При этом нижний индекс около массового числа указывает (в процентах) распространённость изотопа в природе.

Запишите в таблицу число протонов и число нейтронов в ядре самого распространённого изотопа галлия.
\begin{center}
\begin{tabular}{|c|c|}
\hline
Число протонов & Число нейтронов \\
\hline
\rule{0pt}{1.5em} & \\
\hline
\end{tabular}
\end{center}

\begin{center}
\setlength{\tabcolsep}{5pt}
\renewcommand{\arraystretch}{1.15}
\begin{tabular}{|c|c||c|c|c|}
\hline
\(2\) & II &
\begin{tabular}{@{}l@{\hspace{1cm}}r@{}}\(\mathbf{Li}\)&\(3\)\\\multicolumn{2}{c}{\textbf{литий}}\\\multicolumn{2}{c}{\(7_{93}\;6_{7}\)}\end{tabular}&
\begin{tabular}{@{}l@{\hspace{1cm}}r@{}}\(\mathbf{Be}\)&\(4\)\\\multicolumn{2}{c}{\textbf{бериллий}}\\\multicolumn{2}{c}{\(9_{100}\)}\end{tabular}&
\begin{tabular}{@{}l@{\hspace{1cm}}r@{}}\(5\)&\(\mathbf{B}\)\\\multicolumn{2}{c}{\textbf{бор}}\\\multicolumn{2}{c}{\(11_{80}\;10_{20}\)}\end{tabular}\\
\hline
\(3\) & III &
\begin{tabular}{@{}l@{\hspace{1cm}}r@{}}\(\mathbf{Na}\)&\(11\)\\\multicolumn{2}{c}{\textbf{натрий}}\\\multicolumn{2}{c}{\(23_{100}\)}\end{tabular}&
\begin{tabular}{@{}l@{\hspace{1cm}}r@{}}\(\mathbf{Mg}\)&\(12\)\\\multicolumn{2}{c}{\textbf{магний}}\\\multicolumn{2}{c}{\(24_{79}\;26_{11}\;25_{10}\)}\end{tabular}&
\begin{tabular}{@{}l@{\hspace{1cm}}r@{}}\(13\)&\(\mathbf{Al}\)\\\multicolumn{2}{c}{\textbf{алюминий}}\\\multicolumn{2}{c}{\(27_{100}\)}\end{tabular}\\
\hline
\raisebox{-2em}{\(4\)} & IV &
\begin{tabular}{@{}l@{\hspace{1cm}}r@{}}\(\mathbf{K}\)&\(19\)\\\multicolumn{2}{c}{\textbf{калий}}\\\multicolumn{2}{c}{\(39_{93}\;41_{6{,}7}\)}\end{tabular}&
\begin{tabular}{@{}l@{\hspace{1cm}}r@{}}\(\mathbf{Ca}\)&\(20\)\\\multicolumn{2}{c}{\textbf{кальций}}\\\multicolumn{2}{c}{\(40_{97}\;44_{2{,}1}\)}\end{tabular}&
\begin{tabular}{@{}l@{\hspace{1cm}}r@{}}\(\mathbf{Sc}\)&\(21\)\\\multicolumn{2}{c}{\textbf{скандий}}\\\multicolumn{2}{c}{\(45_{100}\)}\end{tabular}\\
\cline{2-5}
 & V &
\begin{tabular}{@{}l@{\hspace{1cm}}r@{}}\(29\)&\(\mathbf{Cu}\)\\\multicolumn{2}{c}{\textbf{медь}}\\\multicolumn{2}{c}{\(63_{69}\;65_{31}\)}\end{tabular}&
\begin{tabular}{@{}l@{\hspace{1cm}}r@{}}\(30\)&\(\mathbf{Zn}\)\\\multicolumn{2}{c}{\textbf{цинк}}\\\multicolumn{2}{c}{\(64_{49}\;66_{28}\;68_{19}\)}\end{tabular}&
\begin{tabular}{@{}l@{\hspace{1cm}}r@{}}\(31\)&\(\mathbf{Ga}\)\\\multicolumn{2}{c}{\textbf{галлий}}\\\multicolumn{2}{c}{\(69_{60}\;71_{40}\)}\end{tabular}\\
\hline
\end{tabular}
\end{center}
}{
\smallskip\noindent\textbf{Решение от GPT}\par
Самый распространённый изотоп галлия имеет массовое число \(A=69\), его распространённость равна \(60\%\). Число протонов равно атомному номеру \(Z=31\), а число нейтронов \(N=A-Z=69-31=38\). Ответ: \(31;\ 38\).
}

% Source BA68F8: https://ege.fipi.ru/bank/questions.php?search=1&proj=BA1F39653304A5B041B656915DC36B38&qid=BA68F8
\BankTask{BA68F8}{
Сколько нейтронов содержится в ядре изотопа платины ${}_{78}^{187}\mathrm{Pt}$ и сколько электронов в электронной оболочке нейтрального атома этого изотопа?

\begin{center}
\begin{tabular}{|c|c|}
\hline
Число нейтронов & Число электронов \\
\hline
\rule{0pt}{1.5em} & \\ \hline
\end{tabular}
\end{center}
}{
\smallskip\noindent\textbf{Решение от GPT}\par
Число нейтронов равно $N=A-Z=187-78=109$. В нейтральном атоме число электронов равно числу протонов: $N_e=Z=78$. Ответ: $109;\ 78$.
}

% Source BFC0DC: https://ege.fipi.ru/bank/questions.php?search=1&proj=BA1F39653304A5B041B656915DC36B38&qid=BFC0DC
\BankTask{BFC0DC}{
Сколько электронов содержится в электронной оболочке нейтрального атома изотопа тория ${}_{90}^{234}\mathrm{Th}$?
}{
\smallskip\noindent\textbf{Решение от GPT}\par
В нейтральном атоме число электронов равно числу протонов: $N_e=Z=90$. Ответ: $90$ электронов.
}

% Source 2B1449: https://ege.fipi.ru/bank/questions.php?search=1&proj=BA1F39653304A5B041B656915DC36B38&qid=2B1449
\BankTask{2B1449}{
Сколько протонов и сколько нейтронов содержится в ядре ${}_{50}^{119}\mathrm{Sn}$?
\begin{center}
\begin{tabular}{|c|c|}
\hline
Число протонов & Число нейтронов \\
\hline
\rule{0pt}{1.5em} & \\
\hline
\end{tabular}
\end{center}
}{
\smallskip\noindent\textbf{Решение от GPT}\par
Число протонов равно зарядовому числу: $Z=50$. Число нейтронов $N=A-Z=119-50=69$. Ответ: $50$ протонов, $69$ нейтронов.
}

\end{document}
